% Part: proof-theory
% Chapter: natural-deduction
% Section: rules-proofs

\documentclass[../../../include/open-logic-section]{subfiles}

\begin{document}

\olfileid{pt}{nat}{rp}

\olsection{નિયમો અને \usetoken{P}{proof}}

પહેલાં કેટલીક વ્યાખ્યાઓ
અને પરિભાષા નક્કી કરીએ.

\Log{N1c} અને~\Log{N1i}માંના
$!A \lif !B$ ધરાવતા
બે નિયમો જુઓ.
\[
\bottomAlignProof
\AxiomC{$\Discharge{!A}{x}$}
\shortDeduce
\DeduceC{$!B$}
\DischargeRule{\Intro{\lif}}{x}
\UnaryInfC{$!A \lif !B$}
\DisplayProof
\qquad
\bottomAlignProof
\AxiomC{$!A \lif !B$}
\AxiomC{$!A$}
\RightLabel{\Elim{\lif}}
\BinaryInfC{$!B$}
\DisplayProof
\]
\Elim{\lif} નિયમ સમજવો
સહેલો છે. રેખાની
ઉપરનાં !!{formula}s
\emph{આધારવાક્યો} છે.
ડાબી બાજુનું !!{formula}
એ \emph{મુખ્ય}
!!{formula}~$!A \lif !B$
છે. દરેક
!!{elimination} નિયમમાં
આવું એક આધારવાક્ય હોય છે,
હંમેશાં સૌથી ડાબે.
તેને અનુમાનનું
\emph{મુખ્ય} આધારવાક્ય
કહે છે. આ કિસ્સાની
જેમ બીજાં આધારવાક્યો
હોય, તો તેમને
\emph{ગૌણ} આધારવાક્યો
કહે છે.
!!{introduction} નિયમોનાં
આધારવાક્યો પણ મુખ્ય
આધારવાક્યો કહેવાય છે.
રેખાની નીચેનું
!!{formula} એ
\emph{ફલિતવાક્ય} છે.

\Intro{\lif} નિયમને
માત્ર એક આધારવાક્ય છે,
અને અહીં ફલિતવાક્ય
એ મુખ્ય સૂત્ર~$!A \lif !B$
છે. બધાં
!!{introduction} નિયમો
આવા છે: ફલિતવાક્ય
મુખ્ય સૂત્ર હોય છે
અને તેનો મુખ્ય સંક્રિયક
``દાખલ'' થાય છે.

સંજ્ઞા~$\Discharge{!A}{x}$નો
અર્થ આ છે. કોઈ
!!a{proof}માં
\Intro{\lif} અનુમાન
આધારવાક્યરૂપે
!!a{formula}~$!B$ને
લાગુ પડે છે.
$!B$માં પૂરી થતી
ઉપ-!!{proof}માં
શરૂઆતનાં કેટલાંક
!!{formula}s હોય છે;
તેમને \emph{ધારણાઓ}
કહે છે.
$!A$ સ્વરૂપની એક
કે વધુ ધારણાઓમાંથી
$!B$ની !!^a{proof}
એ~$!A$માંથી~$!B$
ફલિત થાય છે તેની
!!a{proof} છે.
\Intro{\lif} નિયમ
લગાવીએ ત્યારે
$!A \lif !B$
તારવીએ છીએ,
અને ફલિતવાક્ય
હવે ધારણા~$!A$
પર આધારિત રહેતું નથી.
તે બતાવવા~\Intro{\lif}
અનુમાન ધરાવતી
!!{proof}માંનાં
$!A$ સ્વરૂપનાં
ધારણાઓને \emph{મુક્ત}
અથવા \emph{બંધ}
કહીએ છીએ.
કઈ ધારણાઓ મુક્ત થઈ
શકે તે નિયમ જણાવે છે.
નિષ્કર્ષ~$!A \lif !B$
ધરાવતા~\Intro{\lif}
માટે એ~$!A$ સ્વરૂપની,
એટલે અનુસરણના પૂર્વપદ
સાથે મળતી, ધારણાઓ છે.
કઈ ધારણા કયા નિયમથી
મુક્ત થાય છે તે
બતાવવું ક્યારેક જરૂરી
હોય છે; તે માટે
મુક્તિ-ચિહ્ન~$x$
વાપરીએ છીએ.
પરંતુ ધારણા મુક્ત
કરતો નિયમ જરૂરી
સ્વરૂપની \emph{બધી}
ધારણાઓ, કે એકેય
ધારણા, મુક્ત કરે
તે \emph{જરૂરી નથી}.
દાખલા તરીકે,
પહેલું~\Intro{\lif}
કોઈ ધારણા મુક્ત ન
કરે તોપણ નીચેની
સાબિતી યોગ્ય છે:
\[
\AxiomC{$\Discharge{!A}{y}$}
\RightLabel{\Intro{\lif}}
\UnaryInfC{$!B \lif !A$}
\DischargeRule{\Intro{\lif}}{y}
\UnaryInfC{$!A \lif (!B \lif !A)$}
\DisplayProof
\]
\sourcecorrection{OLMD-231}{મૂળ વૃક્ષના પ્રથમ અનુસરણ-પરિચય પર $x$ મુક્તિ-ચિહ્ન મૂકાયું છે, જ્યારે તે કોઈ ધારણા મુક્ત કરતું નથી અને તરતનું ગદ્ય ચિહ્ન છોડવાનું કહે છે; અહીં તે લેબલ કાઢ્યું છે.}
કોઈ નિયમ એકેય ધારણા
મુક્ત ન કરે ત્યારે
મુક્તિ-ચિહ્ન છોડીએ છીએ
(જેમ અહીં કર્યું છે).

આ !!{proof}માં
\[
\AxiomC{$\Discharge{!A}{x}$}
\AxiomC{$\Discharge{!A}{y}$}
\RightLabel{\Intro{\land}}
\BinaryInfC{$!A \land !A$}
\RightLabel{\Elim{\land}}
\UnaryInfC{$!A$}
\DischargeRule{\Intro{\lif}}{x}
\UnaryInfC{$!A \lif !A$}
\DischargeRule{\Intro{\lif}}{y}
\UnaryInfC{$!A \lif (!A \lif !A)$}
\DisplayProof
\]
પહેલું~\Intro{\lif}
અનુમાન ડાબી
ધારણા~$!A^x$ને,
અને બીજું જમણી
ધારણા~$!A^y$ને
મુક્ત કરે છે.
એટલે~$x$ ચિહ્નવાળી
બધી ધારણાઓને
$x$ ચિહ્નવાળું
\Intro{\lif} મુક્ત કરે છે,
અને~$y$ ચિહ્નવાળીઓને
$y$ ચિહ્નવાળું અનુમાન.
આ રીતે કામ કરવા
એક જ ચિહ્નવાળી
બધી ધારણાઓનું
!!{formula} પણ એક
જ હોવું જરૂરી છે.

પરિમાણકોના નિયમો
થોડા વધુ જટિલ છે.
દાખલા તરીકે,
\Log{N1i}માં
$\lexists$ના નિયમો આ છે:
\[
\bottomAlignProof
\AxiomC{$!A(t)$}
\RightLabel{\Intro{\lexists}}
\UnaryInfC{$\lexists[x][!A(x)]$}
\DisplayProof
\qquad
\bottomAlignProof
\AxiomC{$\lexists[x][!A(x)]$} 
\AxiomC{$\Discharge{!A(c)}{x}$}
\shortDeduce
\DeduceC{$!B$}
\DischargeRule{\Elim{\lexists}}{x}
\BinaryInfC{$!B$} \DisplayProof 
\bottomAlignProof
\]
\Intro{\lexists} નિયમમાં
આધારવાક્ય~$!A(t)$ છે,
જ્યાં~$t$ બંધ પદ છે,
અર્થાત્ ચલ વિનાનું પદ.
આવું અનુમાન ત્યારે યોગ્ય
છે, એટલે કે
\Intro{\lexists}ના
યોજનાત્મક નિયમને
ત્યારે અનુરૂપ છે,
જ્યારે કોઈ
!!{formula}~$!A$,
કોઈ ચલ~$x$
અને કોઈ બંધ પદ~$t$
માટે નિષ્કર્ષ
$\lexists[x][!A]$
અને આધારવાક્ય
$\Subst{!A}{t}{x}$
હોય.
\Elim{\lexists}થી
$!A(c)$ સ્વરૂપની
ધારણાઓ મુક્ત કરી
શકાય છે:
દરેક અનુમાન માટે
કોઈ !!{constant}~$c$
હોય છે અને
$\Subst{!A}{c}{x}$
સ્વરૂપની ધારણાઓ
મુક્ત કરી શકાય છે.
એક અનુમાનમાં
મુક્ત થતી બધી
ધારણાઓમાં એ જ~$c$
વપરાવો જોઈએ.
આ અચળને
\Elim{\lexists}
અનુમાનનો
\emph{સ્વનિર્ધારક ચલ}
કહે છે.
આ અનુમાન માત્ર ત્યારે
યોગ્ય છે જ્યારે~$c$
મુક્ત થતી~$!A(c)$
સિવાયની કોઈ ખુલ્લી
ધારણામાં, $!B$માં
કે~$!A$માં ન આવતો હોય.
આને \emph{સ્વનિર્ધારક
ચલની શરત} કહે છે.

\Log{N1c} અને~\Log{N1i}માંની
!!^{proof}s એ
સૂત્રોનાં સાન્ત વૃક્ષો છે, જે
સ્વયંસિદ્ધો અને
અનુમાન-નિયમોથી
આગમનાત્મક રીતે બને છે.
\sourcecorrection{OLMD-232}{મૂળના આ વાક્યમાં N1-સાબિતીઓને સિક્વન્ટોનાં વૃક્ષો કહેવાય છે, પણ આગળની વ્યાખ્યા અને N1 નિયમ-તકતી પ્રમાણે તે સૂત્રોનાં વૃક્ષો છે; સિક્વન્ટો N2 તંત્રમાં આવે છે.}
દરેક અંતિમ ડાળી
ધારણા છે; તેને
મુક્તિ-ચિહ્ન
લાગેલું હોઈ શકે.
બીજું દરેક
!!{formula}
તેની તરત ઉપરનાં
!!{formula}sમાંથી
તંત્રના કોઈ
અનુમાન-નિયમથી
મળે છે.
વૃક્ષમાં કોઈ
અનુમાનની ઉપર
ધારણા તરીકે આવતા
!!{formula}ને
તેની ઉપરના અનુમાનથી
મુક્ત ન કરાયું હોય,
તો તે અનુમાન સાપેક્ષે
\emph{ખુલ્લું} કહેવાય છે.

\begin{defn}\ollabel{defn:proof}
\Log{N1c} અને~\Log{N1i}માં
\emph{!!^{proof}s} એ
!!{formula}sનાં
સાન્ત વૃક્ષો છે,
જે આગમનાત્મક રીતે
આ પ્રમાણે બને છે:
\begin{enumerate}
  \item\ollabel{defn:prf:assumption}
  $x$ અથવા~$y$ જેવા
  મુક્તિ-ચિહ્નથી ચિહ્નિત
  દરેક !!{formula} પોતે
  !!a{proof} છે.
  માત્ર એક !!{formula}
  ધરાવતી !!a{proof}માં
  એ !!{formula} ખુલ્લી
  ધારણા ગણાય છે.
  \item\ollabel{defn:prf:one-prem}
  જો~$R$ એ એક, બે
  અથવા ત્રણ આધારવાક્યો
  $!A_1$, $!A_2$, $!A_3$
  અને ફલિતવાક્ય~$!A$
  ધરાવતો નિયમ હોય,
  અને~$\delta_1$,
  $\delta_2$, $\delta_3$
  અનુક્રમે~$!A_1$,
  $!A_2$, $!A_3$માં
  પૂરી થતી !!{proof}s
  હોય, તો અનુક્રમે
  નીચેનાં વૃક્ષો
  \[
  \AxiomC{}
  \RightLabel{$\delta_1$}
  \DeduceC{$!A_1$}
  \RightLabel{$R$}
  \UnaryInfC{$!A$}
  \DisplayProof
\qquad
  \AxiomC{}
  \RightLabel{$\delta_1$}
  \DeduceC{$!A_1$}
  \AxiomC{}
  \RightLabel{$\delta_2$}
  \DeduceC{$!A_2$}
  \RightLabel{$R$}
  \BinaryInfC{$!A$}
  \DisplayProof
\qquad
  \AxiomC{}
  \RightLabel{$\delta_1$}
  \DeduceC{$!A_1$}
  \AxiomC{}
  \RightLabel{$\delta_2$}
  \DeduceC{$!A_2$}
  \AxiomC{}
  \RightLabel{$\delta_3$}
  \DeduceC{$!A_3$}
  \RightLabel{$R$}
  \TrinaryInfC{$!A$}
  \DisplayProof
  \]
  !!a{proof} છે, જો
  તેમને રચતી !!{proof}sની
  ધારણાઓનાં
  ચિહ્નો સુસંગત હોય
  (બે જુદાં !!{formula}sને
  એક જ મુક્તિ-ચિહ્ન
  ન અપાયું હોય)
  અને અનુમાનો~$R$નો
  યોગ્ય ઉપયોગ હોય.
  \sourcecorrection{OLMD-233}{મૂળના ત્રણ આધારવાક્યવાળા વૃક્ષમાં $\delta_3$નું નિષ્કર્ષ $!A_2$ લખાયું છે; ઉપર આપેલી અનુક્રમણીકા પ્રમાણે તે $!A_3$ છે.}
  \item નવી !!{proof}માં
  સૌથી ઉપર આવતાં
  !!{formula}નાં સ્થાનો
  $\delta_1$, $\delta_2$
  અથવા~$\delta_3$માં
  ખુલ્લી ધારણા હોય તો
  નવી !!{proof}ની
  \emph{ખુલ્લી ધારણાઓ}
  ગણાય છે, પરંતુ
  નિયમથી મુક્ત થયેલાં
  સ્થાનોને છોડીને.
  જો નિયમ પર~$x$
  (અથવા~$x\, y$)
  ચિહ્ન હોય,
  તો~\Intro{\lif}
  અને~\Intro{\lnot} માટે
  $\delta_1$માંનાં
  $x$ ચિહ્નવાળાં બધાં
  ખુલ્લાં ધારણાસ્થાનો
  મુક્ત થાય છે;
  \Elim{\lexists} માટે
  $\delta_2$માંનાં
  $x$ ચિહ્નવાળાં બધાં,
  અને~\Elim{\lor} માટે
  $\delta_2$માંનાં
  $x$ ચિહ્નવાળાં તથા
  $\delta_3$માંનાં
  $y$ ચિહ્નવાળાં બધાં
  ખુલ્લાં ધારણાસ્થાનો
  મુક્ત થાય છે.
\end{enumerate}
\end{defn}

!!a{proof}ની
આગમનાત્મક રચનામાં
વપરાતી !!{proof}sને
તેની \emph{ઉપ-!!{proof}s}
કહે છે.
અનુમાનનાં આધારવાક્યોમાં
પૂરી થતી
ઉપ-!!{proof}sને
તે અનુમાનની
\emph{તાત્કાલિક
ઉપ-!!{proof}s}
કહે છે.
\Log{N1c} અને~\Log{N1i}
ના નિયમો
\olref{tab:N1}માં છે.
\sourcecorrection{OLMD-234}{મૂળમાં N2c/N2iના નિયમો tab:N1માં છે એમ કહેવાયું છે; નીચે આયાત થતી N1 નિયમ-તકતી અને અહીંની N1 વ્યાખ્યા પ્રમાણે તંત્રોનાં નામ N1c/N1i છે.}

\subfile{rules-N1}

\begin{defn}
!!a{proof}~$\delta$ની
\emph{!!{height}}~$\pheight{\delta}$
આગમનાત્મક રીતે આ
પ્રમાણે વ્યાખ્યાયિત થાય છે.
\begin{enumerate}
  \item જો~$\delta$માં માત્ર
  એક ધારણા હોય, તો
  $\pheight{\delta} = 0$.
  \item જો~$\delta$નો અંત
  એક આધારવાક્યવાળા અનુમાનમાં
  થાય અને તેની તાત્કાલિક
  ઉપ-!!{proof}~$\delta_1$ હોય,
  તો $\pheight{\delta} = \pheight{\delta_1} + 1$.
  \item જો~$\delta$નો અંત
  બે આધારવાક્યવાળા અનુમાનમાં
  થાય અને તેની તાત્કાલિક
  ઉપ-!!{proof}s~$\delta_1$
  અને~$\delta_2$ હોય, તો
  $\pheight{\delta} =
  \max(\pheight{\delta_1}, \pheight{\delta_2}) + 1$.
  \item જો~$\delta$નો અંત
  \Elim{\lor}માં થાય
  અને તેની તાત્કાલિક
  ઉપ-!!{proof}s~$\delta_1$,
  $\delta_2$ અને~$\delta_3$
  હોય, તો
  $\pheight{\delta} =
  \max(\pheight{\delta_1}, \pheight{\delta_2}, \pheight{\delta_3}) + 1$.
\end{enumerate}
જો સિક્વન્ટ~$S$ની
ઊંચાઈ~$\le n$ ધરાવતી
!!a{proof} હોય, તો
$\Proves[n] S$ લખીએ છીએ.
\end{defn}

\begin{ex}
\Log{N1i}માં
કોઈ પણ !!{formula}
પોતે ખુલ્લી ધારણા
હોઈ પોતાની
!!a{proof} ગણાય છે.
એટલે
\[
\Discharge{!A \land !B}{x}
\]
એ~!!a{proof}~$\delta_1$
છે. તેની ઊંચાઈ~$0$ છે.
$\delta_1$માંથી
\Elim{\land}નાં બે
રૂપોમાંથી એકેક લગાવી
બે !!{proof}s~$\delta_2$
અને~$\delta_3$ મળે છે:
\[
\AxiomC{$\Discharge{!A \land !B}{x}$}
\RightLabel{\Elim{\land}[2]}
\UnaryInfC{$!B$}
\DisplayProof
\qquad
\AxiomC{$\Discharge{!A \land !B}{x}$}
\RightLabel{\Elim{\land}[1]}
\UnaryInfC{$!A$}
\DisplayProof
\]
\sourcecorrection{OLMD-235}{મૂળના આરંભિક અને અનુગામી વૃક્ષોમાં $!A \land !B^x$ લખવાથી $x$ માત્ર $!B$ સાથે જોડાય છે; અહીં ધારણા આખા $!A \land !B$ સૂત્રની હોવાથી તેને એક જ મુક્તિ-ચિહ્નથી દર્શાવ્યું છે.}
$\delta_2$ અને~$\delta_3$ના
છેલ્લા અનુમાનની તાત્કાલિક
ઉપ-!!{proof}s સમાન છે:
$!A \land !B$.
બંને !!{proof}sમાં
$!A\land !B$નું
સ્થાન ખુલ્લી ધારણા છે.
આપણે
$\pheight{\delta_2} = \pheight{\delta_3} = 1$
મેળવીએ છીએ.
\sourcecorrection{OLMD-236}{મૂળમાં $\delta_1$ની ઊંચાઈ એક જ વાક્યમાં $0$ અને પછી $1$ કહેવાઈ છે; અહીં એક પગલું ધરાવતી $\delta_2$ તથા $\delta_3$ની ઊંચાઈ $1$ બતાવી છે.}

\Intro{\land} નિયમથી
$!B \land !A$ સ્વરૂપનું
ફલિતવાક્ય મેળવવા
$!B$ અને~$!A$
સ્વરૂપનાં બે આધારવાક્યો
જોઈએ. તેથી નીચેનું
!!a{proof}~$\delta_4$ છે:
\[
\AxiomC{$\Discharge{!A \land !B}{x}$}
\RightLabel{\Elim{\land}[2]}
\UnaryInfC{$!B$}
\LeftSubproofLabel{$\delta_2$}
\AxiomC{$\Discharge{!A \land !B}{x}$}
\RightLabel{\Elim{\land}[1]}
\UnaryInfC{$!A$}
\RightSubproofLabel{$\delta_3$}
\RightLabel{\Intro{\land}}
\BinaryInfC{$!B \land !A$}
\DisplayProof
\]
તેની ઊંચાઈ
$\pheight{\delta_4} = \max(\pheight{\delta_2},
\pheight{\delta_3})+1 = \max(1,1) + 1 = 2$ છે.
તેમાં~$!A \land !B$
ધારણા તરીકે બે વાર
આવે છે અને બંને
સ્થાનો ખુલ્લાં છે.

છેલ્લે~\Intro{\lif}
લગાવી~$\delta_5$
મેળવી શકાય છે:
\[
\AxiomC{$\Discharge{!A \land !B}{x}$}
\RightLabel{\Elim{\land}[2]}
\UnaryInfC{$!B$}
\LeftSubproofLabel{$\delta_2$}
\AxiomC{$\Discharge{!A \land !B}{x}$}
\RightLabel{\Elim{\land}[1]}
\UnaryInfC{$!A$}
\RightSubproofLabel{$\delta_3$}
\RightLabel{\Intro{\land}}
\BinaryInfC{$!B \land !A$}
\DischargeRule{\Intro{\lif}}{x}
\UnaryInfC{$(!A \land !B) \lif (!B \land !A)$}
\DisplayProof
\]
તેની ઊંચાઈ
$\pheight{\delta_4} + 1 = 3$ છે.
\Intro{\lif} અનુમાન
$x$ ચિહ્નવાળી
$!A \land !B$
સ્વરૂપની ધારણાઓનાં
બધાં સ્થાનો મુક્ત કરે છે;
અર્થાત્ તાત્કાલિક
ઉપસાબિતીમાં પહેલાં
ખુલ્લાં રહેલાં બંને
સ્થાનો મુક્ત થાય છે.
$\delta_5$ની એ જ
ધારણાઓ હોવાથી તેમાં
હવે કોઈ ખુલ્લી
ધારણા નથી.
આથી તે પોતાના
અંતિમ-!!{formula}
$(!A \land !B) \lif (!B \land !A)$ની
!!a{proof} છે.
\end{ex}

\end{document}
